\section{Conditioning, clocks, and geometric prediction}
\begin{proposition}[Central density conditioning]\label{prop:conditioning}
Suppose an unweighted and an $f$-weighted joint density satisfy
\[
 n^2p_n(z)=G_n(z)+r_n(z),\qquad
 n^2p_n^f(z)=a_fG_n(z)+r_n^f(z)
\]
on a common central window, with $G_n\ge g_0>0$ and $|r_n|,|r_n^f|\le\varepsilon_n<g_0$. Then their density ratio differs from $a_f$ by at most
\begin{equation}\label{eq:ratio}
          \frac{(1+|a_f|)\varepsilon_n}{g_0-\varepsilon_n}.
\end{equation}
The same bound holds after conditioning on fixed lattice coordinates and any positive-length roof interval contained in that window, including intervals whose lengths shrink with $n$.
\end{proposition}
\begin{proof}
Subtract $a_f$ times the first identity from the second and divide by $G_n+r_n\ge g_0-\varepsilon_n$. For an interval, integrate both identities first; the errors are at most its length times $\varepsilon_n$ and the denominator is at least its length times $g_0-\varepsilon_n$. The length cancels.
\end{proof}
At an exact roof value this is a density-defined version of a regular conditional expectation for almost every value with positive density. It is not conditioning on a positive-probability singleton. It supplies no estimate outside the common central window and no functional path convergence by itself. These are the precise limitations of this interface to stopped-path large deviations.

\begin{proposition}[A conditional physical-clock transfer]\label{prop:clock}
Let $(X_j,\tau_j)\in\R^d\times(0,\infty)$ be stationary dynamical observations with means $(a,\mu)$, $\mu>0$, and a joint functional CLT with covariance $\Sigma$. Assume the uniform law of large numbers for their clock and that initial and terminal residual records on bounded rescaled time intervals are $o_P(\sqrt t)$. Then the reward accumulated up to physical time $t$, centered by $(a/\mu)t$, has limiting covariance per unit time
\begin{equation}\label{eq:clock}
 \frac1\mu Q\Sigma Q^T,\qquad Q=(I_d,-a/\mu).
\end{equation}
For $X=(\kappa_1,\kappa_2,r)$ this transfers induced displacement and collision counts, not the number of experimental bits.
\end{proposition}
\begin{proof}
The clock law gives its inverse index $N(t)/t\to1/\mu$. The centered functional CLT evaluated at this inverse time and the identity $S_{N(t)}\tau=t+o_P(\sqrt t)$ give
\[
 S_{N(t)}X-(a/\mu)t
 = (S_{N(t)}X-N(t)a)-(a/\mu)(S_{N(t)}\tau-N(t)\mu)+o_P(\sqrt t).
\]
The continuous limit permits the random-time substitution. Projection by $Q$ and evaluation at time $1/\mu$ prove \eqref{eq:clock}. The assumed residual estimate includes unfinished induced returns, which are not bounded merely because individual flights are bounded.
\end{proof}

\begin{lemma}[Covariance continuity from a finite-time comparison]\label{lem:covariance}
Suppose a centered correlation matrix satisfies $\|C_j(R)\|\le Ce^{-cj}$ and
$\|C_j(R)-C_j(R')\|\le C\delta^\alpha e^{a j}$, where $\delta=|R-R'|$, $a,c,\alpha>0$. Then the absolutely convergent Green--Kubo covariance satisfies
\[
 \|\Sigma_R-\Sigma_{R'}\|\le C'\delta^{\alpha c/(a+c)}.
\]
If $a=0$, the bound is instead $C'\delta^\alpha(1+|\log\delta|)$.
\end{lemma}
\begin{proof}
Split the defining correlation series at $J=\lfloor\alpha\log(1/\delta)/(a+c)\rfloor$. For $a>0$, the finite sum is bounded by $C\delta^\alpha e^{aJ}$ and the remaining sum by $Ce^{-cJ}$. Both have the stated order. For $a=0$, sum the constant comparison bound over $J+1$ terms and use the exponential tail. Symmetrizing the correlations only changes the constant.
\end{proof}
The hypotheses of this lemma must concern the actual induced observables and their varying domains. Spectral continuity for a different observable norm is not a substitute for this comparison.

\begin{proposition}[Geometric error and the central prediction scale]\label{prop:prediction}
Assume \eqref{eq:LLT} with error $\epsilon_n$, uniformly elliptic covariances, and estimates $(\widehat m,\widehat\Sigma)$ satisfying $|\widehat m-m_R|+\|\widehat\Sigma-\Sigma_R\|\le\omega(\nu)$. If $\sqrt n\,\omega(\nu)\le1$, replacing the Gaussian parameters in \eqref{eq:LLT} changes its right side on a fixed central window by at most $C(1+\sqrt n)\omega(\nu)$. The total prediction error is bounded by this quantity plus $\epsilon_n$.
\end{proposition}
\begin{proof}
The argument of the Gaussian moves by $\sqrt n(\widehat m-m_R)$. On a slightly enlarged compact window the Gaussian and its derivatives in its argument and covariance are bounded uniformly under ellipticity. The mean value theorem gives the bound. Adding the LLT error finishes the proof.
\end{proof}

There is also a direct equilibrium conclusion that needs no local-limit hypothesis. The mean collision rate is
\[
 \lambda(R)=\frac{2R}{\sqrt3/2-\pi R^2}.
\]
Its derivative is $2/A_R+4\pi R^2/A_R^2$, where $A_R=\sqrt3/2-\pi R^2$. On $I$, $A_R>17/100$, using $\sqrt3>173/100$ and $\pi<22/7$, and direct substitution gives $0<\lambda'(R)<110$. Therefore a radius estimate with error at most $\nu$ gives collision-rate error at most $110\nu$. On the circular submodel the average support function identifies the radius, independently of translation. The finite geometric conclusion of \cite{A2G} therefore transfers to this rate with its geometric observation cost, without recovering the entire unknown launch law.

